\chapter{迁移学习与持续学习}
\begin{quote}\small
\textit{本章摘要。}本章讨论模型在数据分布变化后如何继续有效。先区分迁移中的源域、目标域和负迁移，再说明持续学习中的灾难性遗忘、回放、正则化和参数隔离，最后把适配效率、旧任务保持和隐私预算纳入统一评价。主线是：变化来自哪里，哪些知识应该保留，哪些参数可以更新。
\end{quote}
前面各章假定设备和数据环境已经给定，本章首次把“环境会变化”作为显式变量。贯穿案例中的预测提前量、标签语义和安全告警保持不变；变化的是设备、批次、相机或工况的分布，以及模型更新后旧设备是否仍然可靠。
\section{迁移学习}
源域$\mathcal{S}$和目标域$\mathcal{T}$的分布或任务不同。迁移学习利用源域知识改善目标域风险：
\begin{equation}
 R_{\mathcal{T}}(f)=\E_{(x,y)\sim P_\mathcal{T}}\ell(f(x),y).
\end{equation}
预训练--微调冻结部分参数，特征迁移共享表示，域适配则进一步减小域间差异。一个常见的分布对齐目标是最小化源、目标表示的距离，例如最大均值差异（MMD）：
\begin{equation}
 \operatorname{MMD}^2=\left\|\frac1{n_s}\sum_i\varphi(z_i^s)-\frac1{n_t}\sum_j\varphi(z_j^t)\right\|_{\mathcal{H}}^2.
\end{equation}
迁移是否有效取决于源目标相关性。负迁移说明源知识与目标任务冲突，不能把预训练本身视为无条件收益。

\section{持续学习与领域变化}
\paragraph{持续学习。}
持续学习（continual learning）面对按时间到达的任务或数据流，并希望学习新知识而不过度遗忘旧知识。经验回放保存旧样本，参数正则化限制重要参数的变化，参数隔离为不同任务分配子网络。以旧模型$\theta_{old}$为参考的蒸馏正则可写为
\begin{equation}
 \mathcal{L}=\mathcal{L}_{new}+\lambda\,\KL\left(p_{\theta_{old}}(y\mid x)\|p_\theta(y\mid x)\right).
\end{equation}
评估应同时报告新任务性能、旧任务性能和前向/后向迁移，单看最后任务准确率无法识别灾难性遗忘。

\paragraph{领域自适应的两种设定。}
有标签域适配同时使用源域和少量目标标签；无监督域适配只有目标域输入。域对抗训练引入域判别器$D$，编码器试图让$D$无法区分域：
\begin{equation}
 \min_{F,C}\max_D\;\mathcal{L}_{task}(C(F(x_s)),y_s)-\lambda\mathcal{L}_{domain}(D(F(x)),d).
\end{equation}
梯度反转层实现了这一目标的工程化训练，但若源目标语义不对齐，强行域不变可能损害任务判别信息。

\section{参数高效迁移与遗忘控制}
\paragraph{参数高效迁移。}
迁移策略包括冻结骨干、只训练分类头、Adapter、LoRA 和提示学习。冻结更多参数降低过拟合和存储开销，却限制目标域适应能力。少样本目标域应优先进行线性探测和轻量微调，再决定是否解冻深层参数。

\begin{table}[htbp]
\centering
\caption{持续学习方法的知识保持与系统代价}
\scriptsize
\setlength{\tabcolsep}{3pt}
\resizebox{\linewidth}{!}{%
\begin{tabular}{p{2.7cm}p{3.6cm}p{3.7cm}p{3.8cm}}
\hline
方法 & 保持旧知识的机制 & 主要代价 & 更适合的变化情形 \\
\hline
回放 & 直接重看旧样本或核心集 & 存储、隐私和采样偏差 & 有合规缓存预算、旧分布重要 \\
EWC 等正则化 & 约束重要参数偏离旧值 & Fisher 估计误差，任务多时约束累积 & 无法保存原始样本的连续更新 \\
LoRA/Adapter & 将新知识放入低秩或独立模块 & 适配器管理和路由复杂度 & 多设备、多批次的轻量迁移 \\
参数隔离 & 为任务分配独立子空间或子网络 & 模型容量和版本管理增加 & 任务边界清晰、旧任务必须稳定 \\
生成回放 & 用生成器合成旧样本 & 生成偏差可能逐轮累积 & 原始数据不能留存但允许合成 \\
\hline
\end{tabular}}
\end{table}

\paragraph{持续学习的评价量。}
设$A_{i,j}$表示在学习完任务$j$后对任务$i$的性能。平均最终准确率为$\frac1T\sum_iA_{T,i}$，遗忘量可定义为
\begin{equation}
 F=\frac1{T-1}\sum_{i=1}^{T-1}\left(\max_{j\leq T-1}A_{j,i}-A_{T,i}\right).
\end{equation}
任务边界清晰时可使用任务增量学习；类别逐渐出现时属于类增量学习；数据分布连续变化时属于域增量学习。三者的标签空间和推理假设不同，不能只用一个“持续学习”结果概括。

\paragraph{回放预算与隐私。}
回放缓存有限时，可以按类别、时间或代表性选择样本。梯度均值匹配、核心集选择和生成回放试图在小缓存中保留旧知识。若数据不能保存，蒸馏和生成回放可减少原始样本依赖，但需要验证生成器本身是否遗忘。

\paragraph{工业案例：设备域迁移。}
对于不同批次生产的产品，可以把批次作为时间序列任务。实验应依次加入新批次，记录每一批的即时性能和旧批次性能；同时检查负迁移、域对抗训练删除有用信息的风险，以及样本回放与参数正则化在存储和隐私上的代价。

\paragraph{偏移的细分与迁移风险。}
源域和目标域的差异应具体描述为协变量偏移、标签偏移、概念偏移或条件分布变化。协变量偏移下$P(X)$变化而$P(Y\mid X)$近似稳定；标签偏移下类别比例变化；概念偏移下同一观测与标签的关系发生变化。前两类可以尝试重加权或校准，概念偏移则通常需要新的标签和模型更新。

迁移收益取决于源域知识与目标域任务的相关性。若源域特征捕捉的是设备编号、背景纹理或采集协议，微调可能把伪相关带入目标域。实际实验必须与从零训练、冻结骨干、线性探测和全量微调比较，并按目标域样本量绘制学习曲线。

\paragraph{负迁移的识别。}
负迁移不是目标域分数低这么简单，而是迁移方案低于合适的目标域基线。可以比较不同初始化、不同冻结层和不同目标域子群的性能，检查源域表示是否造成优化困难或类别混淆。若目标域只覆盖源域未出现的故障模式，强制共享高层决策边界尤其危险。

\paragraph{EWC 与参数重要性。}
弹性权重固化（EWC）用旧任务参数和 Fisher 信息近似参数重要性：
\begin{equation}
 \mathcal{L}(\theta)=\mathcal{L}_{new}(\theta)+\frac{\lambda}{2}\sum_jF_j(\theta_j-\theta_{old,j})^2.
\end{equation}
重要参数受到更强约束，不重要参数可以适应新任务。它节省原始数据存储，却依赖 Fisher 估计质量；当任务很多时，多个二次惩罚也会积累，最终限制模型学习新知识。

\paragraph{三种增量学习设定。}
任务增量学习在推理时知道任务 ID；类增量学习需要在不断增加的类别中选择，且不能假设任务边界；域增量学习的类别空间不变，但输入分布持续变化。三种设定的难度不同，实验报告必须说明推理阶段是否提供任务 ID，否则相同的准确率无法比较。

\begin{figure}[htbp]
\centering
\begin{tikzpicture}[>=Latex, node distance=0.75cm, every node/.style={font=\small}]
 \node[draw, rounded corners, fill=blue!10, minimum width=2.2cm] (old) {\shortstack{旧知识\\模型$\theta_t$}};
 \node[draw, rounded corners, fill=orange!12, right=of old, minimum width=2.2cm] (stream) {\shortstack{新数据流\\任务$t+1$}};
 \node[draw, rounded corners, fill=green!12, right=of stream, minimum width=2.2cm] (update) {\shortstack{适应机制\\回放/正则}};
 \node[draw, rounded corners, fill=yellow!20, right=of update, minimum width=2.2cm] (new) {\shortstack{更新模型\\$\theta_{t+1}$}};
 \draw[->, thick] (old)--(stream); \draw[->, thick] (stream)--(update); \draw[->, thick] (update)--(new);
 \draw[->, dashed] (new.south) to[bend right=35] node[below]{旧任务评估与遗忘监测} (old.south);
\end{tikzpicture}
\caption{持续学习在新数据到达时更新模型，同时监测旧知识是否被遗忘。}
\end{figure}

\paragraph{回放、正则化与参数隔离。}
回放直接保留旧样本，通常效果稳定但有存储和隐私成本；正则化不保存样本，却可能无法保留多模态旧分布；参数隔离为不同任务分配适配器或子网络，降低干扰但增加模型管理复杂度。生成回放还会引入生成器误差，连续使用时可能发生模型偏差累积。

\paragraph{持续学习的评价闭环。}
除了平均最终准确率，还应报告即时学习、遗忘量、前向迁移、后向迁移、内存占用和更新时间。对工业流数据，还要记录标签延迟和检测到漂移后的响应时间。一个系统如果保持旧任务性能，却需要几天才能更新，仍可能无法满足生产要求。

以视觉质检的新相机连续上线为例，每周到达一个新域，只有少量人工标注。每次更新前保留固定的旧相机测试集、新相机测试集和未知工况测试集；更新后同时测量新域即时性能、旧域遗忘量、未知域拒识率、训练时间和新增存储。回放方案需要报告缓存样本的类别与设备覆盖，EWC 需要报告 Fisher 的估计窗口和正则系数，LoRA/Adapter 需要报告适配器数量、路由规则和合并后推理延迟。

负迁移判别要以“目标域从零或冻结骨干基线”为参照。若新相机的颜色分布变化但缺陷语义稳定，归一化适配和轻量 Adapter 可能足够；若标注规则或产品类别发生变化，强行对齐特征反而会擦除区分信息。连续学习实验还应设置一次回滚：当新模型在旧域或安全关键缺陷上超过风险阈值时，服务回退到上一版，并保留新数据供离线分析，而不是自动覆盖生产模型。

\paragraph{本章小结。}
迁移学习利用已有知识降低新任务成本，持续学习则把适应过程放到时间轴上。二者的共同风险是把偶然相关性当成可迁移知识。可靠方案必须明确域差异、比较无迁移基线、测量负迁移和遗忘，并把更新资源与安全约束纳入评价。